\noindent\textbf{Elección del procedimiento.} De $\psi=Kxy^3$ se obtiene $u=3Kxy^2$, $v=-Ky^3$. La continuidad da $u_x+v_y=3Ky^2-3Ky^2=0$. La aceleración convectiva es $(3K^2xy^4,3K^2y^5)$ y el laplaciano de velocidad $(6Kx,-6Ky)$. En ausencia de gravedad, $\nabla p=-\rho\boldsymbol a+\mu\nabla^2\boldsymbol u$.

\begin{enumerate}

\item Derivar $\psi$ respecto a $y$.

\[\fitmath{u=3\cdot K\cdot x\cdot {y}^{2}}\]

\[\fitmath{u=3\cdot (0{,}01)\cdot (1)\cdot {(1)}^{2}\approx 0{,}03\,\mathrm{m/s}}\]

\item Cambiar el signo de la derivada respecto a $x$.

\[\fitmath{v=-K\cdot {y}^{3}}\]

\[\fitmath{v=-(0{,}01)\cdot {(1)}^{3}\approx -0{,}01\,\mathrm{m/s}}\]

\item Formar $uu_x+vu_y$.

\[\fitmath{\mathrm{ax}=3\cdot {K}^{2}\cdot x\cdot {y}^{4}}\]

\[\fitmath{\mathrm{ax}=3\cdot {(0{,}01)}^{2}\cdot (1)\cdot {(1)}^{4}\approx 0{,}0003\,\mathrm{m/s^{2}}}\]

\item Formar $uv_x+vv_y$.

\[\fitmath{\mathrm{ay}=3\cdot {K}^{2}\cdot {y}^{5}}\]

\[\fitmath{\mathrm{ay}=3\cdot {(0{,}01)}^{2}\cdot {(1)}^{5}\approx 0{,}0003\,\mathrm{m/s^{2}}}\]

\item La fuerza neta por volumen es densidad por aceleración.

\[\fitmath{\mathrm{fx}=\rho\cdot \mathrm{ax}}\]

\[\fitmath{\mathrm{fx}=(900)\cdot (0{,}0003)\approx 0{,}27\,\mathrm{N/m^{3}}}\]

\item Calcular su segunda componente.

\[\fitmath{\mathrm{fy}=\rho\cdot \mathrm{ay}}\]

\[\fitmath{\mathrm{fy}=(900)\cdot (0{,}0003)\approx 0{,}27\,\mathrm{N/m^{3}}}\]

\item Aplicar viscosidad al laplaciano horizontal.

\[\fitmath{\mathrm{vxforce}=\mu\cdot 6\cdot K\cdot x}\]

\[\fitmath{\mathrm{vxforce}=(0{,}5)\cdot 6\cdot (0{,}01)\cdot (1)\approx 0{,}03\,\mathrm{N/m^{3}}}\]

\item El laplaciano vertical tiene signo opuesto.

\[\fitmath{\mathrm{vyforce}=-\mu\cdot 6\cdot K\cdot y}\]

\[\fitmath{\mathrm{vyforce}=-(0{,}5)\cdot 6\cdot (0{,}01)\cdot (1)\approx -0{,}03\,\mathrm{N/m^{3}}}\]

\item Despejar gradiente de presión de Navier--Stokes.

\[\fitmath{\mathrm{px}=-\mathrm{fx}+\mathrm{vxforce}}\]

\[\fitmath{\mathrm{px}=-(0{,}27)+(0{,}03)\approx -0{,}24\,\mathrm{Pa/m}}\]

\item Conservar los signos en la segunda componente.

\[\fitmath{\mathrm{py}=-\mathrm{fy}+\mathrm{vyforce}}\]

\[\fitmath{\mathrm{py}=-(0{,}27)+(-0{,}03)\approx -0{,}3\,\mathrm{Pa/m}}\]

\item Calcular módulo de la fuerza neta.

\[\fitmath{\mathrm{fnet}=\sqrt{{\mathrm{fx}}^{2}+{\mathrm{fy}}^{2}}}\]

\[\fitmath{\mathrm{fnet}=\sqrt{{(0{,}27)}^{2}+{(0{,}27)}^{2}}\approx 0{,}3818377\,\mathrm{N/m^{3}}}\]

\item Calcular módulo de la fuerza viscosa.

\[\fitmath{\mathrm{fvis}=\sqrt{{\mathrm{vxforce}}^{2}+{\mathrm{vyforce}}^{2}}}\]

\[\fitmath{\mathrm{fvis}=\sqrt{{(0{,}03)}^{2}+{(-0{,}03)}^{2}}\approx 0{,}04242641\,\mathrm{N/m^{3}}}\]

\item El módulo de la suma vectorial no es la suma de módulos.

\[\fitmath{\mathrm{gradp}=\sqrt{{\mathrm{px}}^{2}+{\mathrm{py}}^{2}}}\]

\[\fitmath{\mathrm{gradp}=\sqrt{{(-0{,}24)}^{2}+{(-0{,}3)}^{2}}\approx 0{,}3841875\,\mathrm{Pa/m}}\]

\end{enumerate}
